\documentclass{beamer}

\usepackage{cuhkbeamer}


\begin{document}

% ====================== Slide 1: TITLE SLIDE ======================
\begin{frame}
    \title{LAE Market Structure and Competition}
    \author{Chiu Yu KO}
    \institute{Low Altitude Economics}
    \date{Week 2}
  \titlepage
\end{frame}

% ====================== Slide 2: FORMAT & OBJECTIVES ======================
\begin{frame}[t]
  \frametitle{Learning Objectives}

  \begin{itemize}
    \item \textbf{Explain} how market structure shapes competition in different LAE services.
    \item \textbf{Explain} how utilization and route density can reduce average cost.
    \item \textbf{Identify} the main barriers that make entry costly or difficult.
    \item \textbf{Use} simple models to examine price competition and capacity constraints.
    \item \textbf{Evaluate} how Hong Kong and the wider GBA may contribute complementary capabilities.
  \end{itemize}
\end{frame}

% ====================== Slide 3: ICEBREAKER ======================
\begin{frame}[t]
  \frametitle{The Battle for Kowloon}
  
  \textbf{Imagine this happening right now in Hong Kong airspace:}
  \begin{itemize}
      \item \textbf{Operator A} slashes prices on drone cargo flights to the New Territories, flooding the sky with last-mile deliveries.
      \item \textbf{Operator B} counters by securing premium air-taxi slots to the airport, promising executives they will beat the Cross-Harbour Tunnel traffic.
  \end{itemize}

  \vspace{0.3cm}
  \textbf{This is not random chaos. It is a calculated fight for survival.}
  \begin{itemize}
    \item They aren't just fighting for customers. They are fighting for the rarest asset in Hong Kong: \textit{Physical Landing Space}.
    \item \textbf{Quick poll (2 min):} "In 3 years, will one company own the HKIA-Central route, or will there be 10 companies competing?"
  \end{itemize}
\end{frame}

% ====================== Week 2 - Session 1 Title ======================
\begin{frame}[c]
  \begin{center}
    \huge\color{cuhkpurple}\textbf{Session 1}\\
    \vspace{0.5cm}
    \Large Market Segmentation, Concentration \& Economies of Density
  \end{center}
\end{frame}

% ====================== Slide 4: CONCEPT PRIMER 1 ======================
\begin{frame}[t]
  \frametitle{Concept Primer: What is "Market Structure"?}
  
  Before we look at the flying taxis, we need to understand how economists label competition.

  \vspace{0.3cm}
  \textbf{The Three Main Structures:}
  \begin{itemize}
      \item \textbf{Monopoly:} One provider serves a defined market. Its pricing power is still constrained by demand, regulation, and alternative services. (Example: A single toll tunnel).
      \item \textbf{Oligopoly:} A small group of big players (2 to 5 companies). They watch each other carefully. If one drops their price, the others panic. (Example: Mobile phone networks like CSL, 3, SmarTone).
      \item \textbf{Perfect Competition:} Hundreds of small players selling the exact same thing. Nobody has power to raise prices. (Example: Street food stalls).
  \end{itemize}
  
  \vspace{0.3cm}
  \textit{LAE does not have one market structure. Competition may differ by application, corridor, and infrastructure layer.}
\end{frame}

\begin{frame}[t]
\frametitle{Who Might Compete in Hong Kong's LAE Ecosystem?}
\begin{itemize}
\item \textbf{Aircraft developers and manufacturers}
\item \textbf{Transport and logistics operators}
\item \textbf{Telecommunications and digital-platform providers}
\item \textbf{Infrastructure providers}, including landing, charging, and maintenance facilities
\item \textbf{Supporting-service providers}, including finance, leasing, insurance, certification, and professional services
\end{itemize}
\textit{The relevant competitors depend on the specific market being analysed. Sandbox participation alone does not establish market dominance.}
\end{frame}

% ====================== Slide 5: MARKET SEGMENTS ======================
\begin{frame}[t]
  \frametitle{Four Market Segments — Four Different Battles}
  
  Not every low-altitude service competes the same way. The market naturally splits into four clear areas:

  \begin{enumerate}
    \item \textbf{Last-Mile Logistics:} Speed and cost rule here. Drones win when they beat road congestion on predictable B2B delivery routes.
    \item \textbf{Passenger Air Taxis:} Time is the product. Premium travellers pay high prices for minutes saved in the most congested corridors.
    \item \textbf{Industrial Inspection \& Surveying:} Reliability matters most. Companies pay for drones that inspect bridges without shutting down ground traffic.
    \item \textbf{Emergency Medical Services (EMS):} The highest-value segment. Hospitals pay whatever it takes when lives are on the line.
  \end{enumerate}
\end{frame}

% ====================== Slide 6: CONCEPT PRIMER 2 ======================
\begin{frame}[t]
  \frametitle{Concept Primer: ``Market Concentration''}
  
  \textbf{What is Market Concentration?}
  \begin{itemize}
      \item It asks a simple question: \textit{How much of the pie do the biggest players eat?}
      \item We measure this to see if a company has too much power.
  \end{itemize}

  \vspace{0.3cm}
  \textbf{Why does it matter in LAE?}
  \begin{itemize}
      \item If you are the only operator on the Central-HKIA route, you can charge HK\$1,000 per ticket. 
      \item If 5 other operators are allowed to land on your exact same rooftop, you have to drop your price to compete, or you will lose your customers.
  \end{itemize}
\end{frame}

% ====================== Slide 7: STANDALONE FORMULA ======================
\begin{frame}[t, shrink=10]
  \frametitle{Applied Tool: Assessing Market Power (HHI)}
  
  \begin{block}{The Formula (Herfindahl-Hirschman Index)}
  \begin{equation}
      HHI = \sum_{i=1}^{N} s_i^2
  \end{equation}
  \textit{where $s_i$ is the market share (in \%) of firm $i$ in a specific corridor.}
  \end{block}

  \vspace{0.1cm}
  \begin{block}{What This Actually Means for Business in Plain English}
  \begin{itemize}
    \item \textbf{Why do we square the number?} Squaring makes big numbers look much bigger. It mathematically punishes markets where one player is huge. (e.g., $10^2 = 100$, but $90^2 = 8100$).
    \item \textbf{The Scoreboard:} HHI gives a score from near 0 (perfect competition) to 10,000 (a pure monopoly, since $100\% \times 100\% = 10,000$).
    \item \textbf{The Policy Use:} HHI is a screening tool. A high score may justify closer analysis, but does not by itself prove market power or determine a regulatory response.
  \end{itemize}
  \end{block}
\end{frame}

% ====================== Slide 8: BACK-OF-ENVELOPE (HHI) ======================
\begin{frame}[t, shrink=10]
  \frametitle{Back-of-the-Envelope: Do we have a Monopoly?}
 
  \textbf{Scenario:} The government is reviewing market share on the new Northern Metropolis delivery corridor to decide if they need to step in and regulate prices.
  \vspace{0.3cm}
  \begin{columns}[T]
    \begin{column}{0.48\textwidth}
      \begin{block}{Case A: Fragmented Market}
      \begin{itemize}
          \item Four different drone operators share the airspace equally.
          \item Shares: 25\%, 25\%, 25\%, 25\%
          \item $HHI = 25^2 + 25^2 + 25^2 + 25^2$
          \item $HHI = 625 \times 4$
          \item \textbf{HHI = 2,500}
          \item \textit{Result: The market is less concentrated than Case B, although four equal firms still produce an HHI of 2,500.}
      \end{itemize}
      \end{block}
    \end{column}
    \begin{column}{0.48\textwidth}
      \begin{block}{Case B: Sandbox Dominance}
      \begin{itemize}
          \item One early mover locked up the best landing pads before anyone else.
          \item Shares: 80\% (Incumbent), 20\% (Small Rival)
          \item $HHI = 80^2 + 20^2$
          \item $HHI = 6,400 + 400$
          \item \textbf{HHI = 6,800}
          \item \textit{Result: The market is highly concentrated and may warrant closer examination of entry conditions and competitive conduct.}
      \end{itemize}
      \end{block}
    \end{column}
  \end{columns}
\end{frame}

% ====================== Slide 9: CONCEPT PRIMER 3 ======================
\begin{frame}[t]
  \frametitle{Concept Primer: Fixed vs. Marginal Costs}
 
  To understand how flying taxis make money, you must understand two types of costs:
  
  \vspace{0.3cm}
  
  \textbf{1. Fixed Costs (The Rent):}
  \begin{itemize}
      \item Money you have to pay \textit{even if you fly zero customers}.
      \item Example: Leasing a rooftop landing pad in Central costs HK\$1 Million a month, whether you fly 1 time or 1,000 times.
  \end{itemize}
  \vspace{0.3cm}
  \textbf{2. Marginal Costs (The Variable Cost):}
  \begin{itemize}
      \item The tiny extra cost to fly \textit{just one more customer}.
      \item Example: The electricity to charge the drone's battery for one trip across the harbour (maybe HK\$50).
  \end{itemize}
\end{frame}

% ====================== Slide 9.5: TRANSITION ======================
\begin{frame}[t]
  \frametitle{From Concept to Math}
 
  Fixed and marginal costs are the foundation.  
  The next step is to combine them into one powerful number that operators use every day:
  
  \vspace{0.5cm}
  \centering
  \Large\textbf{Average Cost per Flight}
  
  \vspace{0.3cm}
  This tells us the true cost of each flight when we spread the huge fixed costs across many trips.
\end{frame}

% ====================== Slide 10: APPLIED TOOL (ATC) ======================
\begin{frame}[t, shrink=10]
  \frametitle{Applied Tool: Average Cost per Flight}
 
  \begin{block}{The Formula}
  \begin{equation}
      \text{ATC} = \frac{\text{Fixed Costs}}{\text{Number of Flights}} + \text{Marginal Cost}
  \end{equation}
  \end{block}
 
  \vspace{0.1cm}
  \begin{block}{What This Actually Means for Business}
  \begin{itemize}
    \item Fixed Costs (e.g. rooftop rent) stay the same no matter how many times you fly.
    \item The more flights you operate, the lower the Fixed Cost per flight becomes.
    \item Marginal Cost is the extra cost of one additional flight (mainly electricity and minor wear).
  \end{itemize}
  \end{block}
\end{frame}

% ====================== Slide 10.5: BACK-OF-THE-ENVELOPE (ATC) ======================
\begin{frame}[t, shrink=10]
  \frametitle{Back-of-the-Envelope: How Fixed Costs Make Low Utilization Expensive}
 
  \textbf{Scenario:} A rooftop landing pad in Central costs HK\$1,000,000 per month to lease.  
  Marginal cost (electricity + minor maintenance) = HK\$50 per flight.
  \vspace{0.3cm}
  \begin{columns}[T]
    \begin{column}{0.48\textwidth}
      \begin{block}{Low Utilization (20 flights/day)}
      \begin{itemize}
          \item Fixed cost per day ≈ HK\$33,333
          \item Fixed cost per flight: \$33,333 / 20 = \textbf{HK\$1,667}
          \item Marginal cost: HK\$50
          \item \textbf{Total cost per flight: HK\$1,717}
          \item \textit{Ticket price must be very high — only premium customers can afford it.}
      \end{itemize}
      \end{block}
    \end{column}
    \begin{column}{0.48\textwidth}
      \begin{block}{High Utilization (200 flights/day)}
      \begin{itemize}
          \item Fixed cost per day ≈ HK\$33,333
          \item Fixed cost per flight: \$33,333 / 200 = \textbf{HK\$167}
          \item Marginal cost: HK\$50
          \item \textbf{Total cost per flight: HK\$217}
          \item \textit{Ticket price becomes competitive and profitable.}
      \end{itemize}
      \end{block}
    \end{column}
  \end{columns}
 
  \vspace{0.4cm}
  \textit{Takeaway: High fixed costs make low utilization deadly. Operators must maximise flights per day to survive.}

\vfill
{\tiny Illustrative assumptions for classroom analysis; not current market data.}
\end{frame}

% ====================== Slide 11: CONCEPT PRIMER 4 ======================
\begin{frame}[t]
  \frametitle{Concept Primer: The "Density" Secret}
 
  \textbf{Why do traditional airlines want bigger planes?}
  \begin{itemize}
      \item Airlines rely on \textbf{Economies of Scale}. It's cheaper per person to fly one giant Boeing 777 with 300 seats than to fly 15 small planes with 20 seats.
  \end{itemize}
  \vspace{0.3cm}
  \textbf{The LAE Shift: We can't build bigger drones.}
  \begin{itemize}
      \item Flying taxis are stuck at 4 seats. Delivery drones carry 5kg. We can't scale the vehicle.
      \item Instead, we rely on \textbf{Economies of Density}.
      \item Since your Fixed Cost (the rooftop rent) is massive, flying 10 times a day is a financial disaster. Higher route volume can spread fixed costs across more flights, although demand, capacity, congestion, weather, and regulation limit achievable density.
  \end{itemize}
\end{frame}

% ====================== Slide 12: STANDALONE FORMULA ======================
\begin{frame}[t, shrink=10]
  \frametitle{Applied Tool: The Density Cost Advantage}
 
  \begin{block}{The Formula}
  \begin{equation}
      \text{Average Cost per flight (AC)} = \frac{\text{Fixed Route Costs}}{\text{Flights} \times \text{Density}} + \text{Marginal Cost}
  \end{equation}
  \end{block}
  \vspace{0.1cm}
  \begin{block}{What This Actually Means for Business}
  \begin{itemize}
    \item \textbf{Why Density Wins:} If density (the number of flights) is the denominator, then running maximum flights through a single corridor drastically lowers the Average Cost of every single ticket.
    \item \textbf{Potential Incumbent Advantage:} An established operator with greater route volume may sustain a lower average cost than a small entrant.
    \item \textbf{Network Effects:} A broader network may become more useful to customers, although this is conceptually different from economies of density.
  \end{itemize}
  \end{block}
\end{frame}

% ====================== Slide 13: BACK-OF-ENVELOPE (DENSITY) ======================
\begin{frame}[t, shrink=10]
  \frametitle{Back-of-the-Envelope: The Density Trap}
 
  \textbf{Scenario:} Calculating the lowest price we can charge to break even. \\
  \textbf{Assumptions:} A rooftop landing pad (Fixed Cost) is HK\$50,000/day to rent. Electricity/wear (Marginal Cost) is HK\$100/flight.
  \vspace{0.3cm}
  \begin{columns}[T]
    \begin{column}{0.48\textwidth}
      \begin{block}{Operator A: Low Density}
      \begin{itemize}
          \item Volume: Only 20 flights per day.
          \item Pad rent spread out: \$50,000 / 20 = \$2,500
          \item Electricity: \$100
          \item \textbf{Total Cost to Operator: HK\$2,600/flight}
          \item \textit{Result: Operator A must charge HK\$3,000 a ticket just to survive. Only billionaires will fly.}
      \end{itemize}
      \end{block}
    \end{column}
    \begin{column}{0.48\textwidth}
      \begin{block}{Operator B: High Density}
      \begin{itemize}
          \item Volume: 500 flights per day.
          \item Pad rent spread out: \$50,000 / 500 = \$100
          \item Electricity: \$100
          \item \textbf{Total Cost to Operator: HK\$200/flight}
          \item \textit{Result: Operator B has a lower average cost under these assumptions, although feasibility and market price still depend on demand and capacity.}
      \end{itemize}
      \end{block}
    \end{column}
  \end{columns}

\vfill
{\tiny Illustrative assumptions for classroom analysis; not current market data.}
\end{frame}

% ====================== Slide 13: SESSION 2 TITLE ======================
\begin{frame}[c]
  \begin{center}
    \huge\color{cuhkpurple}\textbf{Session 2}\\
    \vspace{0.5cm}
    \Large The Moat: Entry Barriers \& Defense
  \end{center}
\end{frame}

% ====================== Slide 14a: ENTRY BARRIERS - NASA FRAMEWORK ======================
\begin{frame}[t]
  \frametitle{Entry Barriers: The NASA Framework}
 
  \begin{itemize}
    \item In 2018, NASA studied why new drone and flying-taxi companies struggle to enter the market.
    \item They created a simple, practical list of the biggest obstacles.
    \item We still use this list today because it is clear and directly applies to Hong Kong.
    \item These barriers create a huge “moat” around the LAE market.
    \item They make it very difficult and expensive for new startups to start operating.
  \end{itemize}
\end{frame}

% ====================== Slide 14b: THE THREE BARRIERS ======================
\begin{frame}[t]
  \frametitle{The Three Main NASA-Style Entry Barriers}
 
  \begin{enumerate}
    \item \textbf{Regulatory / Certification (The Valley of Death)}  
      Companies must spend millions on lawyers and safety tests for years before they are legally allowed to fly paying customers.
    
    \item \textbf{Infrastructure Availability}  
      Finding physical rooftop space and heavy electrical connections in land-scarce Hong Kong is extremely difficult for a new company.
    
    \item \textbf{Public Perception}  
      A single neighbourhood noise complaint or privacy concern can force the government to shut down a corridor overnight.
  \end{enumerate}
 
  \vspace{0.3cm}
  \textit{The importance of each barrier varies by application and may change as policy, technology, and the market develop.}
\end{frame}


% ====================== Slide 15: CONCEPT PRIMER 5 ======================
\begin{frame}[t]
  \frametitle{Concept Primer: "Sunk Costs"}
  
  \textbf{What is a Sunk Cost?}
  \begin{itemize}
      \item It is money you have already spent that you \textbf{can never get back}, even if you quit the business.
      \item If you buy a van for a delivery business and then quit, you can sell the van to get some money back. That is NOT a sunk cost.
      \item But if you spend HK\$10 Million paying lawyers to get a flight permit, and then you quit... you cannot sell that permit to anyone else. The unrecoverable portion of that expenditure is a sunk cost.
  \end{itemize}

  \vspace{0.3cm}
  \textbf{Why it matters:}
  \begin{itemize}
      \item High Sunk Costs terrify new startups. If a startup knows they have to burn HK\$10M just to enter the market, they will probably stay away. Big companies use this fear to their advantage.
  \end{itemize}
\end{frame}

% ====================== Slide 16: CONCEPT PRIMER 6 ======================
\begin{frame}[t]
  \frametitle{Concept Primer: "Limit Pricing"}
  
  \textbf{How do you defend your turf if you are the first one in?}
  \begin{itemize}
      \item Let's say you are the only flying taxi in Hong Kong. You \textit{could} charge HK\$5,000 a ticket and make huge profits today.
      \item But if your profits are that huge, rich investors in another country will see that, buy their own drones, pay the Sunk Costs, and invade your market.
  \end{itemize}

  \vspace{0.3cm}
  \textbf{The Strategy: Limit Pricing}
  \begin{itemize}
      \item You intentionally do not charge the absolute maximum. 
      \item You charge a price that makes you a *little* bit of money, but is perfectly calculated to guarantee that any new competitor would lose money if they tried to copy you. 
  \end{itemize}
\end{frame}

% ====================== Slide 17: STANDALONE FORMULA ======================
\begin{frame}[t, shrink=10]
  \frametitle{Applied Tool: Limit Pricing (Defending a Route)}
  
  \begin{block}{The Formula}
  \begin{equation}
      P_{\text{Limit}} = AC_{\text{Entrant}} = \frac{\text{Sunk Entry Cost}}{\text{Expected Volume (Q)}} + MC
  \end{equation}
  \end{block}

  \vspace{0.1cm}
  \begin{block}{What This Actually Means for Business}
  \begin{itemize}
    \item \textbf{Weaponising the NASA Barriers:} The NASA barriers (Regulation, Infrastructure) create massive "Sunk Entry Costs" for new rivals. They have to pay millions just to show up to the starting line.
    \item \textbf{The Incumbent's Game:} If you own the route, you charge a price ($P_{\text{Limit}}$) that is exactly equal to the Average Cost ($AC$) of a new entrant.
    \item \textbf{The Core Lesson:} This simplified calculation estimates an entrant’s break-even price. It does not guarantee that limit pricing will prevent entry.
  \end{itemize}
  \end{block}
\end{frame}

% ====================== Slide 18: BACK-OF-ENVELOPE (LIMIT PRICING) ======================
\begin{frame}[t, shrink=10]
  \frametitle{Back-of-the-Envelope: Defending the Corridor}
  
  \textbf{Scenario:} You run a monopoly eVTOL route. A rival wants to enter. They expect to capture 10,000 flights. Their Marginal Cost (MC) is \$200. How do you price your tickets to keep them out?

  \vspace{0.3cm}
  \begin{columns}[T]
    \begin{column}{0.48\textwidth}
      \begin{block}{Case A: High NASA Barriers}
      \begin{itemize}
          \item CAD forces rival to spend years building their own vertiport.
          \item Sunk Entry Cost: HK\$10 Million
          \item $AC_{\text{Entrant}} = \frac{10,000,000}{10,000} + 200$
          \item $AC_{\text{Entrant}} = 1000 + 200 = 1200$
          \item \textbf{Your Limit Price: HK\$1,190}
          \item \textit{Result: You can safely charge high prices. The rival cannot match \$1,190 without losing money.}
      \end{itemize}
      \end{block}
    \end{column}

    \begin{column}{0.48\textwidth}
      \begin{block}{Case B: Low Barriers (Shared Hubs)}
      \begin{itemize}
          \item CAD forces you to share your landing pad with the rival.
          \item Sunk Entry Cost drops to: HK\$1 Million
          \item $AC_{\text{Entrant}} = \frac{1,000,000}{10,000} + 200$
          \item $AC_{\text{Entrant}} = 100 + 200 = 300$
          \item \textbf{Your Limit Price: HK\$290}
          \item \textit{Result: Your moat is gone. You must slash your prices to \$290 to stop the rival from invading.}
      \end{itemize}
      \end{block}
    \end{column}
  \end{columns}

\vfill
{\tiny Illustrative assumptions for classroom analysis; not current market data.}
\end{frame}






% ====================== Slide 19: HK VS CHINA SCALE ======================
\begin{frame}[t]
\frametitle{Hong Kong Experimentation and GBA Scale}
\begin{itemize}
\item \textbf{GBA capabilities:} Strong aircraft, drone, electronics, battery, and supply-chain capabilities.
\item \textbf{Hong Kong's possible contribution:} A complex urban environment together with finance, insurance, logistics, legal, and professional services.
\item \textbf{Sandbox X:} Controlled trials can generate evidence, reduce uncertainty, and refine project concepts and regulatory requirements.
\item \textbf{Strategic opportunity:} Combine complementary capabilities while addressing cross-boundary procedures, customs, data, communications, and safety.
\end{itemize}
\textit{Hong Kong's advantage is a strategy to build, not an automatic outcome.}
\end{frame}

% ====================== Slide 20: SESSION 3 TITLE ======================
\begin{frame}[c]
  \begin{center}
    \huge\color{cuhkpurple}\textbf{Session 3}\\
    \vspace{0.5cm}
    \Large Price Wars \& Co-opetition
  \end{center}
\end{frame}

% ====================== Slide 21: CONCEPT PRIMER 7 ======================
\begin{frame}[t]
  \frametitle{Concept Primer: The "Race to the Bottom"}
  
  \textbf{What happens when two airlines fly the exact same route?}
  \begin{itemize}
      \item If Cathay Pacific and HK Express both have half-empty planes going to Tokyo, what do they do? They drop the price to try and steal the other's customers.
      \item They will keep undercutting each other until the ticket price is exactly equal to the cost of the fuel and the sandwich. All profit disappears.
      \item Economists call this a \textbf{Bertrand Price War}.
  \end{itemize}


\end{frame}

% ====================== Slide 22: STANDALONE FORMULA ======================
\begin{frame}[t, shrink=10]
  \frametitle{Applied Tool: Bertrand Price Competition}
  
  \begin{block}{The Formula}
  \begin{equation}
      P = MC \quad \text{(if capacity is unconstrained)}
  \end{equation}
  \end{block}

  \vspace{0.1cm}
  \begin{block}{What This Actually Means for Business}
  \begin{itemize}
    \item \textbf{The Death of Margins:} If operators have unlimited landing space and empty seats, they will undercut each other until Price ($P$) equals Marginal Cost ($MC$). No one makes a profit.
    \item \textbf{Capacity Saves the Day:} Because HK real estate restricts vertiports to 1 or 2 pads, neither operator can serve the whole market. 
    \item \textbf{The Core Lesson:} Capacity constraints may soften price competition, while reliability, timing, access, and service quality can also differentiate operators.
  \end{itemize}
  \end{block}
\end{frame}

% ====================== Slide 23: BACK-OF-ENVELOPE (BERTRAND) ======================
\begin{frame}[t, shrink=10]
  \frametitle{Back-of-the-Envelope: Capacity Saves the Margin}
  
  \textbf{Scenario:} Operator A and Operator B fly the same route. Electricity/Cost per flight is \$100. Customer Willingness to Pay is \$500. Total Demand is 100 passengers.

  \vspace{0.3cm}
  \begin{columns}[T]
    \begin{column}{0.48\textwidth}
      \begin{block}{Unconstrained Capacity}
      \begin{itemize}
          \item Both operators have huge fleets.
          \item Operator A charges \$400.
          \item Operator B drops price to \$300 to steal all 100 passengers.
          \item Operator A counters at \$200.
          \item \textbf{Final Price drops to \$100 (The cost to fly).}
          \item \textit{Result: A price war destroys the industry. Zero profit.}
      \end{itemize}
      \end{block}
    \end{column}

    \begin{column}{0.48\textwidth}
      \begin{block}{Geofenced / Capacity Constrained}
      \begin{itemize}
          \item Vertiport rules limit both operators to only 30 flights a day each.
          \item Total supply (60) is less than total demand (100).
          \item Operator A charges the max: \$500.
          \item Operator B doesn't need to undercut; they can easily fill their 30 seats at \$500 anyway.
          \item \textbf{Price pressure may be weaker, but the final price need not equal maximum willingness to pay.}
          \item \textit{Result: Limited capacity may support margins, but demand, allocation rules, regulation, and differentiation still matter.}
      \end{itemize}
      \end{block}
    \end{column}
  \end{columns}

\vfill
{\tiny Illustrative assumptions for classroom analysis; not current market data.}
\end{frame}

% ====================== Slide 21: CONCEPT PRIMER 7 ======================



% ====================== NEW SLIDE B: APPLIED TOOL (BERTRAND WITH CONSTRAINT) ======================


% ====================== Back-of-the-Envelope for Slide B ======================



% ====================== Slide 24: CO-OPETITION ======================
\begin{frame}[t]
  \frametitle{The Rise of "Co-opetition" (Frenemies)}

  \textbf{Aviation Benchmark:} Airlines compete for ticket sales, but they share Oneworld lounges and codeshare flights to survive. \\
  \textbf{LAE Shift:} Operators must share fundamental infrastructure to survive the early market phase.

  \begin{itemize}
    \item \textbf{What is Co-opetition?} Competing fiercely on some fronts while cooperating on others to save money and grow the overall market.
    \item \textbf{Shared Vertiports:} Sharing a roof reduces your fixed costs and builds the network faster.
    \item \textbf{Strategic Choice:} You compete on the \textit{vehicle app experience} and brand, but you cooperate on the \textit{landing pads and digital traffic control}.
  \end{itemize}
\end{frame}


% ====================== NEW SLIDE A: CO-OPETITION COST SAVING ======================
\begin{frame}[t, shrink=10]
  \frametitle{Applied Tool: Co-opetition Cost Saving Formula}
 
  \begin{block}{The Formula}
  \begin{equation}
      \text{Savings per Operator} = \frac{\text{Fixed Infrastructure Cost} \times (N-1)}{N}
  \end{equation}
  \end{block}
 
  \vspace{0.1cm}
  \begin{block}{What This Actually Means for Business}
  \begin{itemize}
    \item $N$ = number of operators sharing the same vertiport or UTM system.
    \item Fixed cost (e.g. rooftop rent) is split among more companies.
  \end{itemize}
  \end{block}
\end{frame}

% ====================== Back-of-the-Envelope for Slide A ======================
\begin{frame}[t, shrink=10]
  \frametitle{Back-of-the-Envelope: Co-opetition Savings}
 
  \textbf{Scenario:} A prime Kowloon rooftop vertiport costs HK\$100,000 per day to lease and maintain.
  \vspace{0.3cm}
  \begin{columns}[T]
    \begin{column}{0.48\textwidth}
      \begin{block}{Alone (No Sharing)}
      \begin{itemize}
          \item 1 operator pays full rent.
          \item Fixed cost per day: \textbf{HK\$100,000}
          \item If flying 100 times/day: HK\$1,000 per flight just for rent.
      \end{itemize}
      \end{block}
    \end{column}
    \begin{column}{0.48\textwidth}
      \begin{block}{Sharing with 4 Operators}
      \begin{itemize}
          \item 4 operators split the same rooftop.
          \item Fixed cost per operator: \$100,000 / 4 = \textbf{HK\$25,000}
          \item Savings per operator = HK\$75,000 per day.
          \item If each flies 100 times/day: rent cost drops to HK\$250 per flight.
      \end{itemize}
      \end{block}
    \end{column}
  \end{columns}
  
  \vspace{0.4cm}
  \textit{Takeaway: Sharing infrastructure is one of the fastest ways for early LAE operators to survive Hong Kong’s extreme land costs.}

\vfill
{\tiny Illustrative assumptions for classroom analysis; not current market data.}
\end{frame}

% ====================== Slide 25: GLOBAL COMPARATORS ======================
\begin{frame}[t]
  \frametitle{Global Regulatory Comparators}
  
  How does Hong Kong's strategy compare to the rest of the world?

  \begin{itemize}
    \item \textbf{Hong Kong (Sandbox X):} Controlled trials test project concepts, operating conditions, infrastructure, and cross-boundary procedures.
    \item \textbf{Europe (EU U-space):} Highly structured, digital-first. They are building the software rules before they allow the hardware to fly.
    \item \textbf{USA (FAA Innovate28):} An implementation plan coordinating certification, operations, airspace, infrastructure, environment, and community engagement.
    \item \textbf{Managerial lesson:} Firms must adapt their entry strategies to how each jurisdiction sequences trials, certification, infrastructure, and deployment.
  \end{itemize}
\end{frame}

% ====================== Slide 26: EXERCISE INSTRUCTIONS ======================
\begin{frame}[t]
  \frametitle{In-Class Exercise: Market Simulation}
  
  \textbf{Activity: Oligopoly Pricing Game}
  \begin{itemize}
    \item \textbf{Format:} Groups of 4–5
    \item \textbf{Scenario:} You are one of two operators granted access to a new Northern Metropolis drone delivery corridor.
    \item \textbf{Task:} Using the provided handout:
    \begin{enumerate}
        \item Look at the Payoff Matrix.
        \item Decide if your team will set a "High Price" or a "Low Price".
        \item What happens to your profits if the other team cheats?
    \end{enumerate}
    \item \textbf{Timing:} 15 minutes group work + 5 minute share-out.
  \end{itemize}
\end{frame}

% ====================== Slide 27: EXERCISE TEMPLATE ======================
\begin{frame}[t]
  \frametitle{Exercise Template (Visual)}
  
  \begin{center}
      \textit{[Distribute 1-page printable simulation template]}
      
      \vspace{0.5cm}
      
      \textbf{The Payoff Matrix (Team A vs Team B):} \\
      \vspace{0.2cm}
      If both price High $\rightarrow$ Both survive with healthy margins. \\
      If one prices Low $\rightarrow$ Captures market, but margins evaporate. \\
      \vspace{0.5cm}
      \textit{Discussion: How does a physical capacity constraint (only 1 landing pad) change the math of this game?}
  \end{center}
\end{frame}

% ====================== Slide 28: KEY TAKEAWAYS ======================
\begin{frame}[t]
  \frametitle{Key Takeaways}
  
  \begin{itemize}
    \item LAE market power is influenced by \textbf{infrastructure scarcity} (vertiports) and \textbf{route density}.
    \item NASA barriers to entry are immense; big operators use these massive Sunk Costs to establish \textbf{Limit Pricing}, keeping rivals out.
    \item Capacity constraints and service differentiation may soften price competition, but they do not guarantee high prices or profits.
    \item \textbf{Co-opetition} may reduce duplicated infrastructure cost, but requires careful rules for access, investment, and control.
  \end{itemize}
\end{frame}

% ====================== Slide 29: ASSIGNMENT ======================
\begin{frame}[t]
  \frametitle{Week 2 Assignment Briefing}
  
  \textbf{Two-page Market Competition Memo (max 800 words)}
  \begin{itemize}
    \item \textbf{Due:} Beginning of Week 3
    \item \textbf{Task:} Act as an operator entering a specific HK corridor.
    \item \textbf{Requirements:} 
    \begin{enumerate}
        \item Give a rough score for market concentration (HHI).
        \item Explain how you will achieve "Density" on this specific route.
        \item Classify the top 3 entry barriers (NASA-style) keeping your rivals out.
        \item Recommend one strategic "co-opetition" move to save money.
    \end{enumerate}
    \item \textbf{Grading Focus:} Applying the economic concepts (40\%), realistic HK linkage (30\%), simple math quantification (20\%), professionalism (10\%).
  \end{itemize}
\end{frame}

% ====================== Slide 30: PREVIEW WEEK 3 ======================
\begin{frame}[t]
  \frametitle{Preview of Week 3}
  
  \begin{center}
      \Large \textbf{Cost Structures and Revenue Management}
  \end{center}
  
  \vspace{1cm}
  \textbf{Preview question for next week:} \\
  \textit{How do you stop your business from going bankrupt when a massive Hong Kong rainstorm grounds your entire fleet for 48 hours?}
\end{frame}

% ====================== Slide 31: THANK YOU ======================
\begin{frame}[t]
  \frametitle{Thank You + Q\&A}
  
  \textbf{Pre-Class Readings Reminder (for Week 3):}
  \begin{itemize}
      \item \textit{McKinsey: Unit Economics of eVTOLs}
      \item \textit{GBA Aviation Infrastructure Cost Case Study}
  \end{itemize}
  
  \vspace{1cm}
  \begin{center}
      \textbf{Questions?} \\
      \small (Final 15 minutes reserved for extended Q\&A)
  \end{center}
\end{frame}

\end{document}